% ============================================================================
% LIAL SW01 CHEATSHEET
% Vektoren, Matrizen und lineare Gleichungssysteme
% ============================================================================
% Based on the lecture slides from Hochschule Luzern
% Course: Lineare Algebra (I.BA_LIAL)
% ============================================================================

\documentclass[a4paper, fontsize=7pt]{scrartcl}

% ============================================================================
% ENCODING & LANGUAGE
% ============================================================================
\usepackage[utf8]{inputenc}
\usepackage[german]{babel}

% ============================================================================
% PAGE LAYOUT
% ============================================================================
\usepackage[landscape, left=0.75cm, top=1cm, right=0.75cm, bottom=1.5cm, footskip=5pt]{geometry}
\usepackage{flowfram}
\ffvadjustfalse
\setlength{\columnsep}{1cm}

% Number of columns
\Ncolumn{4}

% Remove paragraph indentation
\setlength{\parindent}{0cm}

% ============================================================================
% COLORS & BOXES
% ============================================================================
\usepackage[most]{tcolorbox}

% Base style for all boxes
\tcbset{
  base/.style={
    boxrule=0.3mm,
    leftrule=1mm,
    left=1.75mm,
    arc=0mm,
    fonttitle=\bfseries,
    colbacktitle=black!10!white,
    coltitle=black,
    toptitle=0.75mm,
    bottomtitle=0.25mm,
    title={#1}
  }
}

% Main concept box (green accent)
\definecolor{maincolor}{rgb}{0.004, 0.5, 0.35}
\newtcolorbox{mainbox}[1]{
  colframe=maincolor,
  base={#1}
}

% Sub concept box (gray accent)
\newtcolorbox{subbox}[1]{
  colframe=black!20!white,
  base={#1}
}

% ============================================================================
% MATHEMATICS
% ============================================================================
\usepackage{amsmath, amsthm, mathtools, amssymb}
\usepackage{amsfonts, amstext}
\allowdisplaybreaks

% Number sets
\def\N{\mathbb{N}}
\def\Z{\mathbb{Z}}
\def\Q{\mathbb{Q}}
\def\R{\mathbb{R}}
\def\C{\mathbb{C}}

% ============================================================================
% TABLES
% ============================================================================
\usepackage{tabularx, multirow, booktabs}
\renewcommand*{\arraystretch}{1.5}

% ============================================================================
% LISTS
% ============================================================================
\usepackage{enumitem}
\setitemize{itemsep=0.5pt, topsep=2pt}
\setenumerate{itemsep=0.75pt, topsep=2pt}

% ============================================================================
% GRAPHICS & TIKZ
% ============================================================================
\usepackage{graphicx}
\usepackage{tikz}
\usetikzlibrary{
  positioning,
  calc,
  arrows.meta
}

\tikzset{
  >=stealth,
}

% ============================================================================
% HYPERLINKS
% ============================================================================
\usepackage{hyperref}
\hypersetup{
  colorlinks=true,
  linkcolor=maincolor,
  urlcolor=maincolor
}

% ============================================================================
% ADDITIONAL PACKAGES
% ============================================================================
\usepackage{centernot}
\usepackage{marvosym}

% ============================================================================
% DOCUMENT METADATA
% ============================================================================
\title{LIAL}
\author{Your Name}
\date{Januar 2026}

\setcounter{tocdepth}{4}
\setcounter{secnumdepth}{4}

% ============================================================================
% DOCUMENT START
% ============================================================================
\begin{document}

\maketitle

% ============================================================================
% SECTION 1: SW01 - VEKTOREN, MATRIZEN, LGS
% ============================================================================
\section{SW01: Vektoren, Matrizen, LGS}

\begin{mainbox}{Kartesisches Koordinatensystem}
Das 2D kartesische Koordinatensystem besteht aus:
\begin{itemize}
    \item \textbf{x-Achse} (Abszisse, horizontal)
    \item \textbf{y-Achse} (Ordinate, vertikal)
\end{itemize}

\textbf{Eigenschaften:}
\begin{itemize}
    \item Achsen stehen senkrecht aufeinander
    \item Gleicher Massstab auf beiden Achsen
    \item Rechtshändiges System (Linksdrehung positiv)
\end{itemize}

\textbf{Punktnotation:} $P(x/y)$ oder $P(x,y)$
\end{mainbox}

\begin{mainbox}{Skalare vs. Vektoren}
\textbf{Skalar:} Nur Betrag (z.B. Temperatur, Druck)

\textbf{Vektor:} Betrag \textbf{und} Richtung
\begin{itemize}
    \item Hat eine \textbf{Länge} (Betrag)
    \item Hat eine \textbf{Richtung}
    \item Geometrisch: Pfeil
    \item Arithmetisch: Zahlen in Klammern
\end{itemize}

Beispiele: Kraft, Geschwindigkeit, Verschiebung
\end{mainbox}

\begin{mainbox}{Vektornotation}
\textbf{Spaltenvektor:}
$$\mathbf{v} = \begin{pmatrix} v_1 \\ v_2 \end{pmatrix} \quad \text{(2D)}, \quad \mathbf{v} = \begin{pmatrix} v_1 \\ v_2 \\ v_3 \end{pmatrix} \quad \text{(3D)}$$

\textbf{Verbindungsvektor:} $\overrightarrow{AB}$ von $A$ nach $B$

\textbf{Ortsvektor:} Vom Ursprung $O$ zum Punkt $P$
$$\overrightarrow{OP} = \begin{pmatrix} p_1 \\ p_2 \end{pmatrix}$$

\textbf{Wichtig:} Zwei Vektoren mit gleicher Länge und Richtung sind \textbf{gleich} (unabhängig vom Anfangspunkt).
\end{mainbox}

\begin{subbox}{Beispiel: Verbindungsvektor}
$A(-7/-3)$, $B(13/-5)$:
$$\overrightarrow{AB} = \begin{pmatrix} 13 - (-7) \\ -5 - (-3) \end{pmatrix} = \begin{pmatrix} 20 \\ -2 \end{pmatrix}$$
\end{subbox}

% ============================================================================
% VEKTOROPERATIONEN
% ============================================================================
\subsubsection{Vektoroperationen}

\begin{mainbox}{Vektoraddition}
\textbf{Geometrisch:} Spitze-an-Fuss-Methode

\textbf{Arithmetisch:}
$$\boxed{\mathbf{a} + \mathbf{b} = \begin{pmatrix} a_1 \\ a_2 \end{pmatrix} + \begin{pmatrix} b_1 \\ b_2 \end{pmatrix} = \begin{pmatrix} a_1 + b_1 \\ a_2 + b_2 \end{pmatrix}}$$

Komponenten werden einzeln addiert.
\end{mainbox}

\begin{subbox}{Beispiel: Addition}
$$\mathbf{a} = \begin{pmatrix} -5 \\ 3 \end{pmatrix}, \quad \mathbf{b} = \begin{pmatrix} 11 \\ 7 \end{pmatrix}$$
$$\mathbf{a} + \mathbf{b} = \begin{pmatrix} -5 + 11 \\ 3 + 7 \end{pmatrix} = \begin{pmatrix} 6 \\ 10 \end{pmatrix}$$
\end{subbox}

\begin{mainbox}{Skalare Multiplikation}
\textbf{Geometrisch:}
\begin{itemize}
    \item $\lambda > 1$: Vektor wird gestreckt
    \item $0 < \lambda < 1$: Vektor wird gestaucht
    \item $\lambda < 0$: Richtung wird umgekehrt
    \item $\lambda = -1$: Gegenvektor
\end{itemize}

\textbf{Arithmetisch:}
$$\boxed{\lambda \mathbf{a} = \lambda \begin{pmatrix} a_1 \\ a_2 \end{pmatrix} = \begin{pmatrix} \lambda a_1 \\ \lambda a_2 \end{pmatrix}}$$
\end{mainbox}

\begin{subbox}{Beispiel: Skalare Multiplikation}
$$\mathbf{a} = \begin{pmatrix} 3 \\ 2 \end{pmatrix}, \quad 3\mathbf{a} = \begin{pmatrix} 9 \\ 6 \end{pmatrix}$$
\end{subbox}

\begin{mainbox}{Vektorsubtraktion}
$$\boxed{\mathbf{a} - \mathbf{b} = \mathbf{a} + (-1) \cdot \mathbf{b} = \begin{pmatrix} a_1 - b_1 \\ a_2 - b_2 \end{pmatrix}}$$

\textbf{Geometrisch:} $\mathbf{a} - \mathbf{b}$ zeigt von der Spitze von $\mathbf{b}$ zur Spitze von $\mathbf{a}$.
\end{mainbox}

\begin{mainbox}{Linearkombination}
Eine \textbf{Linearkombination} der Vektoren $\mathbf{a}, \mathbf{b}, \mathbf{c}$ mit Skalaren $\alpha, \beta, \gamma$:
$$\boxed{\alpha \mathbf{a} + \beta \mathbf{b} + \gamma \mathbf{c}}$$

Allgemein: $\displaystyle\sum_{i=1}^{n} \lambda_i \mathbf{v}_i$
\end{mainbox}

\begin{subbox}{Beispiel: Linearkombination}
$3\mathbf{a} - 2\mathbf{b} + \mathbf{c}$ mit
$\mathbf{a} = \begin{pmatrix} 2 \\ 1 \end{pmatrix}$,
$\mathbf{b} = \begin{pmatrix} -1 \\ 2 \end{pmatrix}$,
$\mathbf{c} = \begin{pmatrix} 1 \\ 3 \end{pmatrix}$:
$$3\begin{pmatrix} 2 \\ 1 \end{pmatrix} - 2\begin{pmatrix} -1 \\ 2 \end{pmatrix} + \begin{pmatrix} 1 \\ 3 \end{pmatrix} = \begin{pmatrix} 9 \\ 2 \end{pmatrix}$$
\end{subbox}

% ============================================================================
% RECHENREGELN
% ============================================================================
\subsubsection{Vektorraumaxiome}

\begin{mainbox}{Rechenregeln für Vektoren}
Für Vektoren $\mathbf{a}, \mathbf{b}, \mathbf{c}$ und Skalare $\lambda, \mu \in \R$:

\renewcommand*{\arraystretch}{1.3}
\begin{tabular}{ll}
$\mathbf{a} + \mathbf{b} = \mathbf{b} + \mathbf{a}$ & Kommutativgesetz \\
$\mathbf{a} + (\mathbf{b} + \mathbf{c}) = (\mathbf{a} + \mathbf{b}) + \mathbf{c}$ & Assoziativgesetz \\
$\mathbf{a} + \mathbf{0} = \mathbf{a}$ & Neutralelement \\
$\mathbf{a} + (-\mathbf{a}) = \mathbf{0}$ & Inverses Element \\
$\lambda(\mathbf{a} + \mathbf{b}) = \lambda\mathbf{a} + \lambda\mathbf{b}$ & Distributivgesetz \\
$(\lambda + \mu)\mathbf{a} = \lambda\mathbf{a} + \mu\mathbf{a}$ & Distributivgesetz \\
$(\lambda\mu)\mathbf{a} = \lambda(\mu\mathbf{a})$ & Assoziativgesetz \\
$1 \cdot \mathbf{a} = \mathbf{a}$ & Neutrales Elem. \\
\end{tabular}

Diese 8 Regeln definieren einen \textbf{Vektorraum}.
\end{mainbox}

% ============================================================================
% SUBSECTION: GERADEN UND EBENEN
% ============================================================================
\subsection{Geraden und Ebenen}

\begin{mainbox}{Punkt-Richtungsform einer Geraden}
Eine Gerade $g$ durch Punkt $P$ (Ortsvektor $\mathbf{p}$) mit Richtungsvektor $\mathbf{u}$:
$$\boxed{g: \mathbf{r} = \mathbf{p} + \lambda \mathbf{u}, \quad \lambda \in \R}$$

$\lambda$ ist der \textbf{Parameter}. Für verschiedene $\lambda$ erhält man verschiedene Punkte auf der Geraden.
\end{mainbox}

\begin{subbox}{Beispiel: Geradengleichung}
$$g: \mathbf{r} = \begin{pmatrix} -3 \\ 2 \end{pmatrix} + \lambda \begin{pmatrix} 2 \\ 2 \end{pmatrix}$$
\begin{itemize}
    \item Stützpunkt: $(-3/2)$
    \item Richtung: $(2/2)$
\end{itemize}
\end{subbox}

\begin{mainbox}{Schnittpunkt zweier Geraden}
Gegeben: $g: \mathbf{r} = \mathbf{p}_1 + \lambda \mathbf{u}$ und $h: \mathbf{r} = \mathbf{p}_2 + \mu \mathbf{v}$

\textbf{Vorgehen:}
\begin{enumerate}
    \item Gleichsetzen: $\mathbf{p}_1 + \lambda \mathbf{u} = \mathbf{p}_2 + \mu \mathbf{v}$
    \item LGS mit 2 Gleichungen und 2 Unbekannten ($\lambda, \mu$) lösen
    \item $\lambda$ oder $\mu$ in Geradengleichung einsetzen
\end{enumerate}

\textbf{Mögliche Lösungen:}
\begin{itemize}
    \item Genau ein Schnittpunkt
    \item Keine Lösung (parallel)
    \item Unendlich viele (identisch)
\end{itemize}
\end{mainbox}

\begin{mainbox}{Punkt-Richtungsform einer Ebene}
Eine Ebene $E$ durch Punkt $P$ (Ortsvektor $\mathbf{p}$) aufgespannt durch $\mathbf{u}$ und $\mathbf{v}$ (nicht parallel):
$$\boxed{E: \mathbf{r} = \mathbf{p} + \lambda \mathbf{u} + \mu \mathbf{v}, \quad \lambda, \mu \in \R}$$

$\mathbf{u}$ und $\mathbf{v}$ heissen \textbf{Spannvektoren}.
\end{mainbox}

\begin{subbox}{Beispiel: Ebenengleichung}
$$E: \mathbf{r} = \begin{pmatrix} 2 \\ 3 \\ 1 \end{pmatrix} + \mu \begin{pmatrix} 4 \\ 2 \\ 9 \end{pmatrix} + \nu \begin{pmatrix} 3 \\ 1 \\ 6 \end{pmatrix}$$
\end{subbox}

\begin{mainbox}{Schnittpunkt Gerade-Ebene}
Gegeben: Gerade $g: \mathbf{r} = \mathbf{p}_g + \lambda \mathbf{u}$ und Ebene $E: \mathbf{r} = \mathbf{p}_E + \mu \mathbf{v} + \nu \mathbf{w}$

\textbf{Vorgehen:}
\begin{enumerate}
    \item Gleichsetzen der Parameterdarstellungen
    \item LGS mit 3 Gleichungen und 3 Unbekannten ($\lambda, \mu, \nu$) lösen
    \item $\lambda$ in Geradengleichung einsetzen
\end{enumerate}

\textbf{Mögliche Lösungen:}
\begin{itemize}
    \item Genau ein Schnittpunkt
    \item Keine Lösung (parallel zur Ebene)
    \item Unendlich viele (Gerade liegt in Ebene)
\end{itemize}
\end{mainbox}

% ============================================================================
% SUBSECTION: MATRIZEN
% ============================================================================
\subsection{Matrizen}

\begin{mainbox}{Definition einer Matrix}
Eine $(m \times n)$-\textbf{Matrix} ist ein rechteckiges Schema von Zahlen mit $m$ Zeilen und $n$ Spalten:
$$\mathbf{A} = \begin{pmatrix} a_{1,1} & a_{1,2} & \cdots & a_{1,n} \\ a_{2,1} & a_{2,2} & \cdots & a_{2,n} \\ \vdots & \vdots & \ddots & \vdots \\ a_{m,1} & a_{m,2} & \cdots & a_{m,n} \end{pmatrix}$$

\textbf{Notation:} $a_{i,j}$ = Element in Zeile $i$, Spalte $j$

\textbf{Kurzschreibweise:} $\mathbf{A} = (a_{i,j})$

\textbf{Quadratische Matrix:} $m = n$
\end{mainbox}

\begin{subbox}{Beispiel: Matrix}
$(2 \times 3)$-Matrix:
$$\mathbf{A} = \begin{pmatrix} 2 & 3 & -1 \\ 1 & -2 & 1 \end{pmatrix}$$
$a_{1,2} = 3$ (1. Zeile, 2. Spalte)
\end{subbox}

\begin{mainbox}{Anwendungen von Matrizen}
\begin{itemize}
    \item Lösen von Gleichungssystemen
    \item Bilder: Graustufenwerte (Matrix), Farbbilder (3 Matrizen für RGB)
    \item Graphentheorie: Adjazenz-/Inzidenzmatrix
    \item PageRank-Algorithmus (Google)
    \item Künstliche neuronale Netze
\end{itemize}
\end{mainbox}

\begin{mainbox}{Matrix-Vektor-Multiplikation}
$(m \times n)$-Matrix $\mathbf{A}$ mal $n$-dimensionaler Vektor $\mathbf{x}$:
$$\boxed{\mathbf{A}\mathbf{x} = \begin{pmatrix} a_{1,1} & \cdots & a_{1,n} \\ \vdots & \ddots & \vdots \\ a_{m,1} & \cdots & a_{m,n} \end{pmatrix} \begin{pmatrix} x_1 \\ \vdots \\ x_n \end{pmatrix} = \begin{pmatrix} \sum_{j=1}^n a_{1,j} x_j \\ \vdots \\ \sum_{j=1}^n a_{m,j} x_j \end{pmatrix}}$$

\textbf{Merke:} Zeile $\cdot$ Spalte $\rightarrow$ Komponente

Ergebnis: $m$-dimensionaler Vektor
\end{mainbox}

\begin{subbox}{Beispiel: Matrix-Vektor-Produkt}
$$\begin{pmatrix} 1 & 2 & 4 \\ 3 & 1 & -1 \end{pmatrix} \begin{pmatrix} x \\ y \\ z \end{pmatrix} = \begin{pmatrix} x + 2y + 4z \\ 3x + y - z \end{pmatrix}$$
\end{subbox}

% ============================================================================
% SUBSECTION: LGS
% ============================================================================
\subsection{Lineare Gleichungssysteme}

\begin{mainbox}{Lineare Gleichung}
Eine \textbf{lineare Gleichung} in einer Variablen $x$:
$$ax + b = 0$$

\textbf{Lösungsfälle:}
\begin{itemize}
    \item $a \neq 0$: Genau eine Lösung $x = -\frac{b}{a}$
    \item $a = 0, b \neq 0$: Keine Lösung
    \item $a = 0, b = 0$: Unendlich viele Lösungen
\end{itemize}
\end{mainbox}

\begin{mainbox}{Lineares Gleichungssystem (LGS)}
\textbf{Allgemeine Form:}
\begin{align*}
a_{1,1}x_1 + a_{1,2}x_2 + \cdots + a_{1,n}x_n &= b_1 \\
a_{2,1}x_1 + a_{2,2}x_2 + \cdots + a_{2,n}x_n &= b_2 \\
\vdots \\
a_{m,1}x_1 + a_{m,2}x_2 + \cdots + a_{m,n}x_n &= b_m
\end{align*}

\textbf{Matrixform:}
$$\boxed{\mathbf{A}\mathbf{x} = \mathbf{b}}$$
\end{mainbox}

\begin{mainbox}{Koeffizientenmatrix}
\textbf{Koeffizientenmatrix} $\mathbf{A}$: Enthält die Koeffizienten $a_{i,j}$

\textbf{Lösungsvektor} $\mathbf{x}$: Enthält die Unbekannten

\textbf{Rechte Seite (RHS)} $\mathbf{b}$: Enthält die Konstanten

$$\underbrace{\begin{pmatrix} a_{1,1} & a_{1,2} \\ a_{2,1} & a_{2,2} \end{pmatrix}}_{\mathbf{A}} \underbrace{\begin{pmatrix} x_1 \\ x_2 \end{pmatrix}}_{\mathbf{x}} = \underbrace{\begin{pmatrix} b_1 \\ b_2 \end{pmatrix}}_{\mathbf{b}}$$
\end{mainbox}

\begin{mainbox}{Erweiterte Koeffizientenmatrix}
Die \textbf{erweiterte Koeffizientenmatrix} enthält $\mathbf{A}$ und $\mathbf{b}$:
$$(\mathbf{A} \mid \mathbf{b}) = \left(\begin{array}{ccc|c} a_{1,1} & \cdots & a_{1,n} & b_1 \\ \vdots & \ddots & \vdots & \vdots \\ a_{m,1} & \cdots & a_{m,n} & b_m \end{array}\right)$$

Diese Form wird für das Gaussverfahren verwendet.
\end{mainbox}

\begin{subbox}{Beispiel: Erweiterte Matrix}
LGS:
$\begin{cases} s - 160t = 0 \\ s - 120t = 100 \end{cases}$

Erweiterte Matrix:
$$\left(\begin{array}{cc|c} 1 & -160 & 0 \\ 1 & -120 & 100 \end{array}\right)$$
\end{subbox}

\begin{mainbox}{Lösungsmethoden}
\begin{itemize}
    \item \textbf{Substitutionsmethode:} Eine Variable isolieren und einsetzen
    \item \textbf{Gleichsetzungsmethode:} Beide Gleichungen nach derselben Variable auflösen
    \item \textbf{Eliminationsmethode (Gauss):} Zeilen addieren/subtrahieren zur Elimination
\end{itemize}

\textbf{Gaussverfahren:} Erweiterte Matrix auf Stufenform bringen, dann \textbf{Rückwärtseinsetzen}.
\end{mainbox}

\begin{subbox}{Beispiel: Elimination}
$$\left(\begin{array}{cc|c} 1 & -160 & 0 \\ 1 & -120 & 100 \end{array}\right) \xrightarrow{Z_2 - Z_1} \left(\begin{array}{cc|c} 1 & -160 & 0 \\ 0 & 40 & 100 \end{array}\right)$$
$\Rightarrow t = 2.5, \quad s = 400$
\end{subbox}

% ============================================================================
% SUBSECTION: SKALARPRODUKT
% ============================================================================
\subsection{Das Skalarprodukt}

\begin{mainbox}{Definition Skalarprodukt}
\textbf{Geometrisch:}
$$\boxed{\mathbf{a} \bullet \mathbf{b} = |\mathbf{a}| \cdot |\mathbf{b}| \cdot \cos \varphi}$$
wobei $\varphi$ der Winkel zwischen $\mathbf{a}$ und $\mathbf{b}$.

\textbf{Arithmetisch (kartesisch):}
$$\boxed{\mathbf{a} \bullet \mathbf{b} = \begin{pmatrix} a_1 \\ a_2 \end{pmatrix} \bullet \begin{pmatrix} b_1 \\ b_2 \end{pmatrix} = a_1 b_1 + a_2 b_2}$$

\textbf{3D:} $\mathbf{a} \bullet \mathbf{b} = a_1 b_1 + a_2 b_2 + a_3 b_3$

\textbf{Ergebnis:} Ein \textbf{Skalar} (Zahl), kein Vektor!
\end{mainbox}

\begin{mainbox}{Betrag (Länge) eines Vektors}
$$\boxed{|\mathbf{a}| = \sqrt{\mathbf{a} \bullet \mathbf{a}} = \sqrt{a_1^2 + a_2^2}}$$

\textbf{3D:} $|\mathbf{a}| = \sqrt{a_1^2 + a_2^2 + a_3^2}$

Dies folgt aus dem Satz des Pythagoras.
\end{mainbox}

\begin{mainbox}{Winkel zwischen Vektoren}
Aus $\mathbf{a} \bullet \mathbf{b} = |\mathbf{a}| \cdot |\mathbf{b}| \cdot \cos \varphi$ folgt:
$$\boxed{\cos \varphi = \frac{\mathbf{a} \bullet \mathbf{b}}{|\mathbf{a}| \cdot |\mathbf{b}|}}$$

$$\boxed{\varphi = \arccos\left(\frac{\mathbf{a} \bullet \mathbf{b}}{|\mathbf{a}| \cdot |\mathbf{b}|}\right)}$$
\end{mainbox}

\begin{mainbox}{Einheitsvektor}
Ein \textbf{Einheitsvektor} hat Länge 1.

Normierung von $\mathbf{a}$ auf Länge 1:
$$\boxed{\hat{\mathbf{a}} = \frac{\mathbf{a}}{|\mathbf{a}|} = \frac{1}{|\mathbf{a}|} \mathbf{a}}$$

Es gilt: $|\hat{\mathbf{a}}| = 1$
\end{mainbox}

\begin{mainbox}{Orthogonalität}
Zwei Vektoren $\mathbf{a}$ und $\mathbf{b}$ sind \textbf{orthogonal} (senkrecht), wenn:
$$\boxed{\mathbf{a} \bullet \mathbf{b} = 0 \quad \Leftrightarrow \quad \mathbf{a} \perp \mathbf{b}}$$

(Der Nullvektor ist zu jedem Vektor orthogonal.)
\end{mainbox}

\begin{subbox}{Beispiel: Orthogonaler Vektor}
Finde $\mathbf{b}$ orthogonal zu $\mathbf{a} = \begin{pmatrix} 2 \\ 1 \end{pmatrix}$:

$\mathbf{a} \bullet \mathbf{b} = 2b_1 + b_2 = 0$

$\Rightarrow b_2 = -2b_1$

Z.B. $\mathbf{b} = \begin{pmatrix} 1 \\ -2 \end{pmatrix}$ oder $\mathbf{b} = \begin{pmatrix} -1 \\ 2 \end{pmatrix}$
\end{subbox}

\begin{mainbox}{Rechenregeln Skalarprodukt}
Für $\lambda \in \R$ und Vektoren $\mathbf{a}, \mathbf{b}, \mathbf{c}$:

\renewcommand*{\arraystretch}{1.3}
\begin{tabular}{ll}
$\mathbf{a} \bullet \mathbf{b} = \mathbf{b} \bullet \mathbf{a}$ & Kommutativ \\
$\mathbf{a} \bullet (\mathbf{b} + \mathbf{c}) = \mathbf{a} \bullet \mathbf{b} + \mathbf{a} \bullet \mathbf{c}$ & Distributiv \\
$\lambda(\mathbf{a} \bullet \mathbf{b}) = (\lambda\mathbf{a}) \bullet \mathbf{b} = \mathbf{a} \bullet (\lambda\mathbf{b})$ & \\
\end{tabular}

\textbf{Hinweis:} $\bullet$ und $\cdot$ werden oft gleich geschrieben.
\end{mainbox}

% ============================================================================
% SUBSECTION: SCHNELLREFERENZ
% ============================================================================
\subsection{Schnellreferenz}

\begin{mainbox}{Wichtige Formeln}
\renewcommand*{\arraystretch}{1.5}
\begin{center}
\begin{tabularx}{\linewidth}{lX}
\textbf{Addition:} & $\mathbf{a} + \mathbf{b} = (a_1+b_1, a_2+b_2)^T$ \\
\textbf{Skalarmult.:} & $\lambda\mathbf{a} = (\lambda a_1, \lambda a_2)^T$ \\
\textbf{Betrag:} & $|\mathbf{a}| = \sqrt{a_1^2 + a_2^2}$ \\
\textbf{Skalarprodukt:} & $\mathbf{a} \bullet \mathbf{b} = a_1 b_1 + a_2 b_2$ \\
\textbf{Winkel:} & $\cos\varphi = \frac{\mathbf{a} \bullet \mathbf{b}}{|\mathbf{a}||\mathbf{b}|}$ \\
\textbf{Orthogonal:} & $\mathbf{a} \perp \mathbf{b} \Leftrightarrow \mathbf{a} \bullet \mathbf{b} = 0$ \\
\textbf{Einheitsv.:} & $\hat{\mathbf{a}} = \frac{\mathbf{a}}{|\mathbf{a}|}$ \\
\textbf{Gerade:} & $\mathbf{r} = \mathbf{p} + \lambda\mathbf{u}$ \\
\textbf{Ebene:} & $\mathbf{r} = \mathbf{p} + \lambda\mathbf{u} + \mu\mathbf{v}$ \\
\textbf{LGS:} & $\mathbf{A}\mathbf{x} = \mathbf{b}$ \\
\end{tabularx}
\end{center}
\end{mainbox}

\begin{subbox}{Dimensionen n-D}
Alle Formeln gelten analog für $n$-dimensionale Vektoren:
$$\mathbf{a} \bullet \mathbf{b} = \sum_{i=1}^{n} a_i b_i, \quad |\mathbf{a}| = \sqrt{\sum_{i=1}^{n} a_i^2}$$
\end{subbox}

\end{document}